% -----------------------------------------------------------------------------------------------------
\section{喷注、喷注子结构与非全局观测量}\label{sec:front-jets}

带非全局对数的 $e^+e^-$ 喷注过程的因子化定理\cite{Becher:2015hka,Becher:2016mmh}——硬函数 $\mathcal H_m$ 乘多 Wilson 线软函数 $\mathcal S_m$，重整化群方程是多重数空间中的矩阵，领头对数与大 $N_c$ 下给出 BMS 方程，窄锥时出现 coft 模式——是核心陈述\cref{thm:jets-ngl-factorization}（[A]）；轻喷注质量与半球软函数\cite{Becher:2016omr}亦然。本节记录超出它的结果；在强子对撞中 Glauber 相位产生超领头对数（\cref{sec:front-sll}）。

\begin{proposition}[非全局对数：超出领头对数与测量驱动的因子化 \tagA]\label{prop:front-ngl}
(a) 色密度矩阵在一切对数阶重求和非全局对数，其反常维度已知到 LO 与 NLO；经共形变换它等价于激波（Regge）演化，重新导出 NLO BFKL／Balitsky--JIMWLK，并在一切对数阶证明发散的指数化（\cite{Caron-Huot:2015bja}，据摘要）。(b) 在大 $N_c$ 与次领头对数精度上，非全局对数的 NLL 重求和：平面大 $N_c$ QCD 中的积分—微分方程，以维数正规化形式与显式有限的四维形式给出，并在 $\order{\alpha_s^2}$ 与完整 QCD 核对\cite{Banfi:2021owj}；喷注演算\cite{Banfi:2021xzn}；两圈反常维度\cite{Becher:2021urs}；以大 $N_c$ 部分子簇射实现的 LHC 上的 NLL 重求和\cite{Becher:2023vrh}；2026 年关于非全局与聚类对数的两圈反常维度\cite{Becher:2026zon}以及 Higgs 产生中 NLL 非全局对数\cite{Becher:2026iim}的预印本（均据摘要）。(c) 通过足够多的测量（“修饰胶子”），非全局对数可以因子化并重求和到一个未分辨的标度（\cite{Larkoski:2015zka}，据摘要）。（本命题只涉及部分子截面。）
% errtype: E-S
\uses{thm:jets-ngl-factorization, thm:jets-bms}
\end{proposition}

\begin{proofsketch}
(a) 把软辐射写成作用于 Wilson 线色密度矩阵的演化，经共形映射与激波演化对应；(b) 在大 $N_c$ 下推导并数值求解 NLL 的演化方程，与固定阶核对；(c) 增加测量把非全局结构化为全局的。完整论证见所引文献。
\end{proofsketch}

\paragraph{地位与缺口。}[A]（(b) 限于大 $N_c$ 与 NLL；只读到摘要层次）。第十六章的\cref{prop:jets-ngl-beyond}（[P]）涉及有限 $N_c$ 与固定阶结果，与这里限于大 $N_c$ 的 [A] 并不矛盾。

\begin{proposition}[喷注子结构、修饰喷注与半单举喷注函数 \tagA]\label{prop:front-jet-substructure}
(a) 在 soft-drop 幂次计数 $z_{\rm cut}\ll1$ 中，一大类 soft-drop 修饰的观测量（含喷注质量）因子化；修饰掉软宽角辐射后，一切奇异贡献都是共线的，因此归一化分布在所有阶都没有非全局对数，并且除夸克／胶子份额外与过程无关；NNLL 重求和\cite{Frye:2016aiz}。(b) 修饰喷注中的强子化：在中间的喷注质量区，领头幂次修正是三个普适的非微扰参数乘以微扰系数；更小的质量处是一个未归一化的形状函数\cite{Hoang:2019ceu}。(c) 半单举喷注函数 $J_i(z,\omega_J,\mu)$ 精确地服从类时 DGLAP 演化（\cref{prop:pdf-timelike-dglap}；$\ln R$ 的重求和），单举喷注产生与单举强子产生共享硬函数，只需 $D_i^h\to J_i$\cite{Kang:2016mcy}。（均据摘要。本命题只涉及部分子截面；(b) 的非微扰参数作为输入。）
% errtype: E-S
\uses{def:jets-algorithms, prop:pdf-timelike-dglap}
\end{proposition}

\begin{proofsketch}
在 SCET 中按 soft-drop 的幂次计数分离硬、全局软、共线—软与共线模式并逐一因子化；(c) 由喷注函数的单圈计算与重整化群相容性得到。完整论证见所引文献。
\end{proofsketch}

\paragraph{地位与缺口。}[A]（(c) 为 NLO 加重整化群相容性；只读到摘要层次）。

\begin{proposition}[能量关联子的共线因子化与共形场论中的光线 OPE \tagA]\label{prop:front-lightray-ope}
(a) 能量—能量关联子在一般量子场论的共线极限下领头幂次因子化，反常维度由类时劈裂函数的矩给出（\cref{prop:pdf-timelike-dglap}）；QCD 中到 NNLL（\cite{Dixon:2019uzg}，据摘要）。(b) 在共形场论中（非微扰），同一类光平面上的零积分（光线）算符有收敛的 OPE（“光线 OPE”；事件形状的天球块），在非微扰共形场论中用 Lorentz 反演公式导出（\cite{Kologlu:2019mfz}，据摘要）——这是事件形状共线因子化的非微扰类比，而不是 QCD 陈述。（(a) 是微扰陈述；(b) 是共形场论中的非微扰陈述。）
% errtype: E-S
\uses{def:jets-algorithms, prop:pdf-timelike-dglap}
\end{proposition}

\begin{proofsketch}
(a) 用光线算符的 OPE 把能量关联子的共线极限写成扭度 2 光线算符之和，其演化由类时劈裂函数的矩控制；(b) 在共形场论中用 Lorentz 反演公式构造光线算符乘积的收敛展开。完整论证见所引文献。
\end{proofsketch}

\paragraph{地位与缺口。}[A]；(b) 是共形场论中的非微扰结论，也是本章唯一的非微扰正面结论，它不涉及 QCD。
